\documentclass[11pt]{article}

\usepackage{amsmath,amssymb,amsthm,mathtools}
\usepackage[T1]{fontenc}
\usepackage{microtype}
\usepackage{hyperref}
\usepackage{enumitem}
\usepackage{booktabs}

\title{Delayed Finite-Window Visibility and Non-Rigidity of\\
Power-Root One-Relator Pro-\texorpdfstring{$p$}{p} Groups}
\author{Seo Seongkyo}
\date{October 5, 2026}

\newtheorem{theorem}{Theorem}[section]
\newtheorem{proposition}[theorem]{Proposition}
\newtheorem{lemma}[theorem]{Lemma}
\newtheorem{corollary}[theorem]{Corollary}
\newtheorem{remark}[theorem]{Remark}
\newtheorem{definition}[theorem]{Definition}
\newtheorem{example}[theorem]{Example}

\newcommand{\Fp}{\mathbb F_p}
\newcommand{\Zp}{\mathbb Z_p}
\newcommand{\D}{D}
\newcommand{\F}{F}
\newcommand{\Dn}{D_n}
\newcommand{\W}{W}
\newcommand{\Wn}{W_n}
\newcommand{\Gal}{\operatorname{Gal}}
\newcommand{\Aut}{\operatorname{Aut}}
\newcommand{\ordz}{\operatorname{ord}_{Z}}
\newcommand{\ab}{\mathrm{ab}}

\begin{document}
\maketitle

\begin{abstract}
We study finite-window visibility and filtration compression for pro-$p$ groups with hidden power-root relations.  Our first result is a universal subgroup-depth comparison: if $K$ is an open subgroup of index $p^s$ in an arbitrary pro-$p$ group $G$, then
\[
D_n(G)\cap K\subseteq D_{\lceil n/p^s\rceil}(K)
\]
for every $n$.  We then prove that the factor $p^s$ is uniformly optimal for every $n$ in the free pro-$p$ class, using explicit commutators of Schreier generators rather than an abstract graded lift.

We apply this filtration theorem to a quadratic one-relator stress family
\[
G_{s,a}=\left\langle z,x_1,\ldots,x_d\ \middle|\
z^{p^s}=x_1^{p^a}r_D\right\rangle,
\qquad
r_D=[x_1,x_2][x_3,x_4]\cdots[x_{d-1},x_d],
\]
with odd $p$, even $d$, and nondegenerate alternating quadratic relation.  The hidden root depth is invisible through depth $p^s$, while the intrinsic transfer obstruction survives at the critical window $p^s+1$.  We obtain the exact finite-window separation threshold
\[
n_{\mathrm{sep}}(s)=p^s+1
\]
in the declared stress-family scope.

Ordinary mod-$p$ cohomology and the quadratic initial relation are blind to the hidden depth.  We also show that the stronger degree-only statement for arbitrary defining relations is false, using $r=z^p$ as a boundary counterexample.

The paper therefore combines a universal filtration theorem, an every-$n$ sharpness theorem in the free-pro-$p$ class, and an exact critical-window theorem for the designated stress family.  The distinction between relative marked information and genuinely intrinsic finite-window information is kept explicit throughout.  At the critical window, the finite parameters $a<s$ are separated already by abelianization, the range $a>s$ is saturated with the power-free case, and the remaining boundary $a=s$ versus $a=\infty$ is separated by the intrinsic transfer obstruction.  No theorem is asserted for arbitrary
defining relations merely from their Zassenhaus order.
\end{abstract}

\section{Introduction}

A recurring question in the preceding finite-window work is whether a
parameter hidden in a pro-$p$ presentation can be recovered from a finite
quotient
\[
\Wn(G)=G/\Dn(G).
\]
The basic model is a root relation
\[
z^{p^s}=r,
\]
where the right-hand side is fixed while the exponent $p^s$ varies.
The first phenomenon is elementary but fundamental: the relation
\(z^{p^s}\) has Zassenhaus depth \(p^s\).  Consequently the parameter
cannot be seen before the window reaches depth \(p^s+1\).

The second phenomenon is subtler.  At the critical depth one must
distinguish three questions:
\begin{enumerate}[label=(\roman*)]
\item does the root relation survive in the finite window?
\item does it survive in a \emph{marked relative} extension over a fixed
quotient?
\item can the resulting information be recovered from the abstract
finite group after the marking is forgotten?
\end{enumerate}
The first two questions can be answered in the stress family considered
here.  The third is a different reconstruction problem and remains open
in the unresolved cases.

The purpose of this paper is therefore not to claim a universal
classification theorem.  Its contribution is a controlled finite-window
theorem together with a sharp description of what the finite window does
and does not remember.

The paper has four main results.

\begin{enumerate}
\item A universal delayed-window lemma shows that for
\(n\le p^s\),
\[
\Wn(G_{s,r})\cong \F/(\Dn(\F),r),
\]
independently of $s$.
\item For the standard quadratic stress family, the marked relative
extension has exact critical threshold \(p^s+1\).
\item The same family is non-rigid with respect to several coarse
invariants: its abelianization, quadratic initial relation, and ordinary
mod-$p$ cohomology do not recover $s$.
\item The assumption \(\ordz(r)\ge2\) alone is insufficient for a
general root-visibility theorem.  Thus the quadratic stress theorem is
a genuine scoped theorem, not merely the first case of a degree-only
universal statement.
\end{enumerate}

\section{Notation and the finite-window problem}

Fix an odd prime \(p\).  Let \(\F\) be a finitely generated free pro-$p$
group and let \(\Dn(\F)\) denote its Zassenhaus filtration.  For a
pro-$p$ group \(G\),
\[
\Dn(G)=\prod_{ij\ge n}\gamma_i(G)^{p^j},
\qquad
\Wn(G)=G/\Dn(G).
\]
We use the standard functoriality
\[
\varphi(\Dn(G))\subseteq \Dn(H)
\]
for every continuous homomorphism \(\varphi:G\to H\).

For a word \(r\in\F\), define
\[
G_{s,r}=\F/\overline{\langle\!\langle z^{p^s}r^{-1}\rangle\!\rangle}.
\]
When
\[
r=r_D=[x_1,x_2][x_3,x_4]\cdots[x_{d-1},x_d],
\]
we use the stress family
\[
G_{s,a}
=
\left\langle z,x_1,\ldots,x_d
\ \middle|\
z^{p^s}=x_1^{p^a}r_D
\right\rangle ,
\qquad a\ge1,
\]
and
\[
G_{s,\infty}
=
\left\langle z,x_1,\ldots,x_d
\ \middle|\
z^{p^s}=r_D
\right\rangle .
\]

The map
\[
\pi_{s,a}:G_{s,a}\longrightarrow D_a
\]
to the Demushkin quotient obtained by killing the $z$-direction will be
called the \emph{marked relative quotient map}.  The exact form of
\(D_a\) is not needed for the delayed-window lemma; it matters only for
the relative critical theorem.

\begin{definition}[Relative finite-window separation]
Let
\[
\pi_n:\Wn(G)\longrightarrow Q_n
\]
be a finite-window quotient map induced by a fixed marked quotient
\(\pi:G\to D\).  We say that relative separation occurs at depth \(n\)
if the corresponding filtered extension over \(Q_n=D/\Dn(D)\) is
non-split, whereas it is split at all smaller depths in the range under
consideration.
\end{definition}

\begin{remark}
Relative separation is a marked notion.  It remembers the specified
map to \(D/\Dn(D)\).  It is therefore strictly weaker than a theorem
saying that the abstract finite groups \(\Wn(G)\) are non-isomorphic.
This distinction will be maintained throughout.
\end{remark}

\section{Universal delayed visibility}

The first result does not use the Demushkin relation.

\begin{lemma}[Delayed root visibility]
Let \(r\in\F\), and let
\[
G_{s,r}=\F/
\overline{\langle\!\langle z^{p^s}r^{-1}\rangle\!\rangle}.
\]
If \(n\le p^s\), then
\[
\boxed{
\Wn(G_{s,r})
\cong
\F/\bigl(\Dn(\F),r\bigr).
}
\]
In particular, for fixed \(r\), the windows of all parameters \(t\ge s\)
are canonically identified for every \(n\le p^s\).
\end{lemma}

\begin{proof}
Since \(z\in\Dn(\F)\) only at depth $1$ and the $p^s$-power has
Zassenhaus weight $p^s$,
\[
z^{p^s}\in\Dn(\F)
\qquad\text{for }n\le p^s.
\]
Hence the relation
\(z^{p^s}r^{-1}\) reduces to \(r^{-1}\) in
\(\F/\Dn(\F)\), giving the displayed quotient.  The construction is
natural in the presentation and does not depend on $s$ once
\(n\le p^s\).
\end{proof}

\begin{corollary}[Universal lower bound]
Any finite-window procedure that is required to distinguish the root
depth \(s\) from larger root depths using only the quotient
\(\Wn(G_{s,r})\) must have
\[
n\ge p^s+1.
\]
This is a lower bound only; it does not assert that the critical window
actually separates the corresponding abstract finite groups.
\end{corollary}

\section{Universal subgroup-depth compression}

The filtration comparison used below is valid without freeness.

\begin{theorem}[Universal Zassenhaus subgroup-depth comparison]
Let $G$ be an arbitrary pro-$p$ group and let $K\triangleleft G$ be an open subgroup with
\[
[G:K]=p^s.
\]
Then for every $n\ge1$,
\[
\boxed{
D_n(G)\cap K\subseteq D_{\lceil n/p^s\rceil}(K).
}
\tag{SC$_s$}
\]
\end{theorem}

\begin{proof}
We first treat $[G:K]=p$.  Choose $a\in G$ with $G/K=\langle aK\rangle$ and put
\[
A=\mathbb F_p[[K]],\qquad B=\mathbb F_p[[G]],\qquad
J=I_K,\qquad t=a-1.
\]
As a left $A$-module,
\[
B=\bigoplus_{r=0}^{p-1}At^r.
\]
Since $a^p\in K$ and the coefficient field has characteristic $p$,
\[
t^p=a^p-1\in J.
\]
Normality gives $aJa^{-1}=J$.

For $m\ge1$ define the weighted normal-form filtration
\[
E_m=
\bigoplus_{r=0}^{p-1}
J^{\max(0,\lceil(m-r)/p\rceil)}t^r,
\]
with the explicit convention $J^0=A$.  A normal-form monomial
$ct^r$, $c\in J^q$, has weight $pq+r$.  If
$\sigma(c)=aca^{-1}$, then
\[
tc=\sigma(c)t+(\sigma(c)-c),
\]
and both coefficients on the right remain in $J^q$.  Together with
$t^p\in J$, this shows that normal-form multiplication does not decrease
weight.  Hence
\[
E_mE_\ell\subseteq E_{m+\ell}.
\]
Since $I_G=JB+tB\subseteq E_1$, we obtain
\[
I_G^n\subseteq E_n.
\]
Taking the direct $A$-component gives
\[
I_G^n\cap A\subseteq J^{\lceil n/p\rceil}.
\]
Using the standard Jennings--Lazard dimension-subgroup identity
\[
D_n(H)=H\cap(1+I_H^n)
\]
for pro-$p$ groups \cite{DDMS1999} proves
\[
D_n(G)\cap K\subseteq D_{\lceil n/p\rceil}(K).
\]

For $[G:K]=p^s$, the finite $p$-group $G/K$ has a composition series.
Pulling it back gives
\[
G=K_0>K_1>\cdots>K_s=K,
\qquad [K_{i-1}:K_i]=p.
\]
Iterating the one-step result and using
\[
\left\lceil\frac{\lceil m/p\rceil}{p}\right\rceil
=
\left\lceil\frac m{p^2}\right\rceil
\]
gives the stated $p^s$-step comparison.
\end{proof}

\begin{remark}
An earlier proof route asserted a stronger augmentation-ideal equality
and is false.  The theorem above uses the weighted normal-form filtration
instead; no such equality is required.
\end{remark}

\section{Every-$n$ sharpness in the free pro-$p$ class}

The factor $p^s$ in $(\mathrm{SC}_s)$ cannot be uniformly improved, and this
optimality holds for every integer $n$, not only for $p$-power values.

Let
\[
F=\langle a,b\rangle
\]
be free pro-$p$, and let
\[
K=\ker\bigl(F\to\mathbb Z/p^s\bigr),
\qquad a\mapsto1,\quad b\mapsto0.
\]
Then $K$ has Schreier free basis
\[
c_0=a^{p^s},\qquad
c_i=a^iba^{-i}\quad(0\le i<p^s).
\]
Write
\[
L_F=\operatorname{gr}_D(F),\qquad
L_K=\operatorname{gr}_D(K).
\]
These are the corresponding free restricted Lie algebras.

For $c_0$,
\[
\operatorname{in}_F(c_0)=\xi^s(\bar a),
\qquad \nu_F(c_0)=p^s,
\]
while
\[
\operatorname{in}_F(c_i)
=(\operatorname{ad}\bar a)^i(\bar b),
\qquad \nu_F(c_i)=i+1,
\]
and each $c_i$ has $K$-degree one.

\begin{proposition}[Every-$n$ sharpness]
For every $n\ge1$ there exists $g_n\in K$ such that
\[
\boxed{
\nu_F(g_n)=n,
\qquad
\nu_K(g_n)=\left\lceil\frac n{p^s}\right\rceil.
}
\]
Consequently
\[
D_n(F)\cap K\not\subseteq
D_{\lceil n/p^s\rceil+1}(K)
\]
for every $n$.
\end{proposition}

\begin{proof}
Write
\[
n=p^sq+r,\qquad 0\le r<p^s.
\]

If $r>0$, define
\[
g_{q,r}
=
[\underbrace{c_0,[c_0,\ldots,[c_0}_{q\text{ times}},c_{r-1}]\ldots]].
\]
The $K$-initial term is
\[
(\operatorname{ad}\bar c_0)^q(\bar c_{r-1}),
\]
which is nonzero because the Schreier generators are free generators of
$L_K$.  Hence
\[
\nu_K(g_{q,r})=q+1.
\]
On the $F$ side the initial term is
\[
(\operatorname{ad}\xi^s(\bar a))^q
(\operatorname{ad}\bar a)^{r-1}(\bar b).
\]
To see that this is nonzero, embed the free restricted Lie algebra in
the free associative algebra $\mathbb F_p\langle a,b\rangle$.
The distinguished word in
$(\operatorname{ad}a)^{r-1}(b)$ is $a^{r-1}b$, while
$\xi^s(a)=a^{p^s}$.  Repeated commutation with $a^{p^s}$ produces the
distinguished word
\[
a^{qp^s+r-1}b
\]
with coefficient $1$, so cancellation is impossible.  Therefore
\[
\nu_F(g_{q,r})=qp^s+r=n.
\]

If $r=0$ and $q\ge1$, one generator of $F$-degree $p^s$ is not enough:
its self-bracket vanishes.  The second such generator is
\[
u_*=
\operatorname{in}_F(c_{p^s-1})
=
(\operatorname{ad}\bar a)^{p^s-1}(\bar b),
\qquad \nu_F(u_*)=p^s.
\]
Set
\[
h_2=[c_0,c_{p^s-1}],
\qquad
h_{k+1}=[h_k,c_0].
\]
The $K$-initial term of $h_k$ is a nonzero iterated bracket of two
distinct free generators, so $\nu_K(h_k)=k$.  On the $F$ side its
distinguished associative word is
\[
a^{(k+1)p^s-1}b
\]
with coefficient $1$, obtained first from
\[
[a^{p^s},a^{p^s-1}b]
=
a^{2p^s-1}b-a^{p^s-1}ba^{p^s}
\]
and then by repeated commutation with $a^{p^s}$.  Thus
\[
\nu_F(h_k)=kp^s.
\]
Taking $g_n=h_q$ completes the case $r=0$.

These are actual group commutators of Schreier generators.  No
injectivity or simultaneous lifting assertion for
$\operatorname{gr}_D(K)\to\operatorname{gr}_D(F)$ is required.
\end{proof}

\begin{corollary}[Uniform optimality]
For every $n\ge1$,
\[
\boxed{
\left\lceil\frac n{p^s}\right\rceil
\text{ is best possible in the free-pro-}p\text{ class}.
}
\]
No corresponding pointwise sharpness claim is made for arbitrary
pro-$p$ groups.
\end{corollary}

\section{The quadratic stress family}

We now specialize to
\[
r_D=[x_1,x_2][x_3,x_4]\cdots[x_{d-1},x_d].
\]
Its image in the degree-two Zassenhaus quotient is the nonzero quadratic
form
\[
\rho_D=[X_1,X_2]+[X_3,X_4]+\cdots+[X_{d-1},X_d].
\]
The power terms \(z^{p^s}\) and \(x_1^{p^a}\) have Zassenhaus degree at
least \(p\ge3\).  Thus the quadratic initial relation is independent of
$s$ and $a$.

\begin{proposition}[Coarse invariance of the quadratic jet]
For fixed odd \(p\) and fixed rank \(d\), all groups \(G_{s,a}\) in the
declared stress family have the same degree-two Zassenhaus initial
relation \(\rho_D\).  In particular, the parameter \(s\) is not visible
in the quadratic initial form.
\end{proposition}

\begin{proof}
For odd \(p\), every $p$-power of a degree-one generator lies in the
third Zassenhaus term.  Hence the degree-two component of the relator is
exactly the commutator product \(r_D\), independent of both power
parameters.
\end{proof}

\section{Intrinsic critical-window separation}

The universal comparison is now combined with an intrinsic transfer
construction.  Let
\[
W=W_{p^s+1}(G)
\]
in the declared quadratic stress-family scope.  The nondegenerate
alternating quadratic initial form determines a canonical one-dimensional
radical direction in $H^1(W,\Fp)$ and hence a canonical index-$p$
kernel
\[
K=\ker(\chi).
\]
Thus no presentation-dependent choice of $z$ is used in defining the
following predicate.

Let
\[
S_s(W):=
\operatorname{im}\bigl(
W^{\mathrm{ab}}[p^s]\longrightarrow
W^{\mathrm{ab}}/pW^{\mathrm{ab}}
\bigr).
\]
In the declared stress-family scope this is one-dimensional over
$\Fp$.  Choose any
$t\in W^{\mathrm{ab}}[p^s]$ whose image spans $S_s(W)$, and let
\[
V:W^{\mathrm{ab}}\longrightarrow K^{\mathrm{ab}}
\]
be the canonical transfer.  Define
\[
\boxed{
\varepsilon_s(W):=
p^{s-1}V(t)
\pmod{p^sK^{\mathrm{ab}}}
\in K^{\mathrm{ab}}/p^sK^{\mathrm{ab}}.
}
\tag{$\varepsilon_s$}
\]
This is independent of the chosen generator of the line up to a unit,
and changing the lift within the same line changes the displayed class
by an element of $p^sK^{\mathrm{ab}}$.  Hence its zero/nonzero value is
an intrinsic invariant of the finite group $W$.

The filtration comparison gives the required truncation control.  At
the critical depth,
\[
D_{p^s+1}(F)\cap K
\subseteq
D_{p^{s-1}+1}(K),
\]
and the defining Zassenhaus product shows
\[
\operatorname{im}
\bigl(D_{p^{s-1}+1}(K)\to K^{\mathrm{ab}}\bigr)
\subseteq p^sK^{\mathrm{ab}},
\]
because all commutator factors disappear after abelianization and the
first surviving pure-power exponent is $p^s$.  Thus
\[
\boxed{
\operatorname{im}
\bigl(D_{p^s+1}(F)\cap K\to K^{\mathrm{ab}}\bigr)
\subseteq p^sK^{\mathrm{ab}}.
}
\tag{TF$_s$}
\]

For the boundary $a=s$, let $A_i$ be the conjugate Schreier generators
for $x_1$ and let $U=z^p$ in the index-$p$ kernel, with
$\sigma(A_i)=A_{i+1}$.  The abelianized Schreier relation induced by
\[
z^{p^s}=x_1^{p^s}r_D
\]
is
\[
p^{s-1}U-p^sA_i=0.
\]
Since $r_D\in[K,K]$, it contributes nothing in $K^{\mathrm{ab}}$.
The critical transfer term is represented by
\[
p^{s-1}(\sigma-1)^{p-1}A_0.
\]
Modulo the relation lattice and modulo $p^s$, the relation
$p^{s-1}U=0$ leaves
\[
(\sigma-1)^{p-1}A_0
=
\sum_{j=0}^{p-1}
(-1)^{p-1-j}\binom{p-1}{j}A_j,
\]
whose first and last coefficients are units modulo $p$.  Therefore
\[
\boxed{
\varepsilon_s(W_{p^s+1}(G_{s,s}))\ne0.
}
\]

For $a=\infty$, the power term $x_1^{p^a}$ is absent.  Taking the
normalized torsion-line representative $t=z$ gives
\[
p^{s-1}V(t)=0
\pmod{p^sK^{\mathrm{ab}}},
\]
so
\[
\boxed{
\varepsilon_s(W_{p^s+1}(G_{s,\infty}))=0.
}
\]
Consequently
\[
\boxed{
W_{p^s+1}(G_{s,s})
\not\cong
W_{p^s+1}(G_{s,\infty}).
}
\]

Together with the lower-window blindness, this yields
\[
\boxed{
n_{\mathrm{sep}}(s)=p^s+1
}
\]
for the declared odd-$p$, even-rank, nondegenerate quadratic stress
family.  The old order-jump route is not used in this conclusion.

\section{The critical relative theorem}

The critical depth is
\[
n_s=p^s+1.
\]

\begin{theorem}[Exact marked relative threshold]
For the declared odd-$p$ stress family, the marked relative finite-window
extension is split for every
\[
n\le p^s
\]
and is non-split at
\[
n=p^s+1.
\]
Equivalently,
\[
\boxed{
n_{\mathrm{sep}}^{\mathrm{rel}}(s)=p^s+1.
}
\]
This statement includes the finite parameters \(a\ge1\) and the
power-free case \(a=\infty\) in the declared stress-family scope.
\end{theorem}

\begin{proof}[Proof architecture]
The proof has two independent parts.

\emph{Lower bound.}
For \(n\le p^s\), Lemma 3.1 identifies the finite window with the
quotient in which the root term \(z^{p^s}\) has disappeared.  The
generator-lift map therefore gives a section of the marked relative
quotient.

\emph{Critical survival.}
At \(n=p^s+1\), the root term can no longer be discarded.  For
\(a\ge2\), an explicit metabelian quotient detects the surviving
\(p^s\)-power and the associated integral Fox divisibility obstruction.
For \(a=1\), a finite semidirect-product pushout gives an independent
witness:
\[
A=\langle x,z\mid z^{p^{s+1}}=1,\;x^p=z^{p^s}\rangle
\]
with an element $y$ acting by
\[
yzy^{-1}=z^{1-p},
\qquad
yxy^{-1}=x.
\]
The relation \(x^p[x,y]=z^{p^s}\) then gives a quotient of
\(G_{s,1}\).  Jennings' formula shows
\[
\Dn(H_s)=1
\qquad(n=p^s+1)
\]
for the relevant finite witness, while \(z^{p^s}\ne1\).  The Fox
derivative calculation
\[
f_x=N_p(x)+x^p-y,\qquad f_y=x^{p+1}-1
\]
reduces in the pure $Y$-direction to
\[
(p-Y)A(Y)=p^s.
\]
Its formal solution has coefficient \(p^{-1}\) in degree $s$, so no
integral section change kills the residual class.  Hence the marked
extension is non-split.

The case \(a=\infty\) is the corresponding power-free stress quotient.
The same critical-depth argument is the base root case.  Thus the
lower and upper bounds agree at \(p^s+1\).
\end{proof}

\begin{remark}
The relative theorem is the extension-theoretic form of the critical
phenomenon.  In the final form of the paper, the marked obstruction is
refined by the intrinsic transfer predicate below, which also closes
the boundary $G_{s,s}$ versus $G_{s,\infty}$ in the declared scope.
\end{remark}

\section{Stress-family non-rigidity}

The delayed-visibility theorem would be much less informative if the
parameter $s$ were already encoded in elementary invariants.  The next
results show that this is not the case.

\subsection{Abelianization}

\begin{proposition}[Abelianization is blind to the root depth]
For finite \(a\),
\[
G_{s,a}^{\ab}
\cong
\mathbb Z_p^d\oplus\mathbb Z/p^a\mathbb Z,
\]
independently of \(s\).
At the own-critical window
\[
n_s=p^s+1,
\]
one has
\[
\boxed{
\W_{n_s}(G_{s,a})^{\ab}
\cong
(\mathbb Z/p^{s+1}\mathbb Z)^d\oplus\mathbb Z/p^a\mathbb Z.
}
\]
\end{proposition}

\begin{proof}
Abelianization replaces the defining relation by
\[
p^s z=p^a x_1.
\]
Smith normal form gives
\[
\Zp^{d+1}/\langle p^s e_z-p^a e_1\rangle
\cong
\mathbb Z_p^d\oplus\mathbb Z/p^a\mathbb Z
\]
when \(a<s\).  In the finite window \(n_s=p^s+1\), the abelian
Zassenhaus filtration contributes \(p^{s+1}\)-torsion in the free
directions, giving the displayed finite abelianization.
\end{proof}

\begin{remark}
The own-critical abelianization separates the sequence of critical
windows indexed by their \emph{own} depths.  It does not solve the
same-numerical-window comparison with a group whose own critical depth
lies later.
\end{remark}

\subsection{Complete critical-window classification of $a$}
\begin{proposition}[Critical-window $a$-classification]
For fixed $s\ge2$ in the declared stress-family scope, the critical
windows satisfy the following trichotomy:
\begin{enumerate}
\item if $1\le a<s$, then
\[
\W_{p^s+1}(G_{s,a})^{\ab}
\cong
(\mathbb Z/p^{s+1}\mathbb Z)^d\oplus\mathbb Z/p^a\mathbb Z,
\]
so distinct values of $a<s$ are separated by abelianization;
\item if $a>s$, then
\[
\W_{p^s+1}(G_{s,a})\cong
\W_{p^s+1}(G_{s,\infty});
\]
\item if $a=s$, then the abelianizations agree with the power-free
case, but
\[
\varepsilon_s(\W_{p^s+1}(G_{s,s}))\ne0,
\qquad
\varepsilon_s(\W_{p^s+1}(G_{s,\infty}))=0,
\]
so the two critical windows are non-isomorphic.
\end{enumerate}
Thus the only nontrivial critical boundary after abelianization is
$a=s$ versus $a=\infty$.
\end{proposition}

\begin{proof}
The first statement is the Smith-normal-form calculation in the preceding
abelianization proposition.  If $a>s$, then $p^a\ge p^{s+1}$, so the
term $x_1^{p^a}$ vanishes in the quotient at depth $p^s+1$, giving the
second statement.  When $a=s$, the same abelianization as the power-free
case occurs, while the intrinsic transfer calculation above gives the
displayed nonzero/zero values of $\varepsilon_s$.  Hence the final
boundary is separated.
\end{proof}

\subsection{\texorpdfstring{Ordinary mod-$p$ cohomology}{Ordinary mod-p cohomology}}

\begin{theorem}[Ordinary cohomology blindness]
Within the declared odd-$p$ stress family, including finite
\(a\ge1\) and \(a=\infty\),
\[
\boxed{
H^\bullet(G_{s,a},\Fp)
\cong
H^\bullet(G_{t,b},\Fp)
}
\]
as graded \(\Fp\)-algebras for all allowed parameter choices.

More explicitly,
\[
H^0\cong\Fp,\qquad
H^1\cong\Fp^{d+1},\qquad
H^2\cong\Fp,\qquad
H^k=0\quad(k\ge3),
\]
and the cup product in degree one is determined by the common quadratic
form \(\rho_D\).
\end{theorem}

\begin{proof}
For odd $p$, the power terms \(z^{p^s}\) and \(x_1^{p^a}\) lie in the
third Zassenhaus term.  Hence every relator in the family has the same
quadratic commutator initial form
\[
[x_1,x_2][x_3,x_4]\cdots[x_{d-1},x_d].
\]
The standard one-relator cohomology result for this quadratic-commutator
congruence gives a quadratic mod-$p$ cohomology algebra
\cite{Quadrelli2022}, concentrated
in degrees at most two, with degree-one cup product determined by the
displayed commutator form.  Thus no $s$- or $a$-dependence remains in
ordinary graded cohomology.
\end{proof}

\begin{remark}
The cohomology theorem is a blindness theorem, not an isomorphism
theorem for the groups.  In particular it says nothing about higher
Bockstein data, Massey products, filtered extension data, or the
abstract finite-window isomorphism problem.
\end{remark}

\section{A failed degree-only generalization}

It is tempting to replace the quadratic stress hypothesis by the single
condition
\[
\ordz(r)\ge2.
\]
That generalization is not justified.

\begin{proposition}[The Zassenhaus order alone is insufficient]
There is no theorem of the form
\[
\ordz(r)\ge2
\quad\Longrightarrow\quad
\text{the full relation }z^{p^s}=r
\text{ is visible at }p^s+1
\]
without an additional critical-survival hypothesis.

\end{proposition}

\begin{proof}
Take $r=z^p$, which has Zassenhaus order $p$.  Then
\[
z^{p^s}=z^p
\]
simplifies to
\[
z^{p^s-p}=1.
\]
The actual relator therefore has a lower $p$-power scale than the
nominal root depth $p^s$.  Its first nontrivial filtered contribution is
already controlled by the degree-$p$ term $z^p$.  Hence the critical
threshold cannot be inferred from the single datum
$\ordz(r)=p$.  This gives the obstruction to the proposed
degree-only arbitrary-$r$ theorem.
\end{proof}
\begin{remark}
This counterexample does not disprove all possible arbitrary-$r$
theorems.  It shows that the correct hypothesis must control the
critical survival of the right-hand side in a non-tautological way.
In particular, a degree-only theorem is not the right generalization.
\end{remark}

\section{What the finite window does not yet reconstruct}

The preceding results leave a precise boundary.

\subsection{Relative versus unmarked information}

The marked theorem controls a diagram
\[
\W_{p^s+1}(G_{s,a})
\longrightarrow
D/\Dn(D).
\]
The abstract finite group
\(\W_{p^s+1}(G_{s,a})\), by contrast, comes without the distinguished
quotient map, the generator $z$, or a chosen orientation character.

A section change in a central Zassenhaus extension changes a cocycle by
a coboundary.  Consequently a scalar defect extracted from one chosen
lift is not automatically an intrinsic invariant of the abstract
finite group.  The natural object after quotienting this gauge freedom
is the extension-equivalence class itself, which is essentially the
structured finite-window isomorphism problem.

This explains why the successful marked theorem cannot simply be
relabelled as an unmarked reconstruction theorem.

\section{Scope, novelty boundary, and future generalization}

The theorem proved here has three deliberate scope restrictions.

\begin{enumerate}
\item The exact critical threshold is a theorem for the declared
stress family, including the intrinsic boundary $a=s$ versus
$a=\infty$ described above.
\item The ordinary cohomology blindness result concerns the full
graded ring \(H^\bullet(-,\Fp)\), not higher operations or filtered
integral data.
\item No claim is made that an arbitrary relation \(r\) with
\(\ordz(r)\ge2\) satisfies the same critical theorem.
\end{enumerate}

The natural next generalization is therefore not another degree-only
search.  A meaningful broader theorem would need a filtered initial
or lift invariant that controls critical survival of \(r\).  Any such
theorem must be functorial under changes of presentation and invariant
under the relevant gauge action.

The present paper intentionally stops before that point.  This makes
the logical status of the result explicit:
\[
\boxed{
\text{quadratic stress-family critical-window theorem}
\quad+\quad
\text{coarse non-rigidity}
\quad+\quad
\text{explicit boundary of the arbitrary-relation generalization}.
}
\]

\section{Conclusion}

Paper 4 establishes a two-level finite-window principle.

First, for every arbitrary pro-$p$ group and every open subgroup of index
$p^s$,
\[
D_n(G)\cap K\subseteq D_{\lceil n/p^s\rceil}(K).
\]
Second, this bound is not merely qualitative: in the free pro-$p$ class
the factor $p^s$ is uniformly optimal for every $n$, with explicit
Schreier-generator commutator witnesses satisfying
\[
\nu_F(g_n)=n,
\qquad
\nu_K(g_n)=\left\lceil\frac n{p^s}\right\rceil.
\]

Applied to the designated quadratic stress family, the same filtration
control closes the intrinsic transfer obstruction and yields the exact
critical window
\[
\boxed{n_{\mathrm{sep}}(s)=p^s+1.}
\]
Thus the hidden root depth is universally invisible through depth $p^s$
but becomes detectable at the next layer in the declared family.

The paper also establishes a sharp limitation: the Zassenhaus degree of
an arbitrary right-hand side $r$ alone does not force the same
critical-visibility theorem.  The counterexample $r=z^p$ rules out that
degree-only generalization.

The resulting contribution is therefore not a universal reconstruction
theorem for arbitrary one-relator pro-$p$ groups.  It is a precise
combination of universal filtration infrastructure, every-$n$ sharpness
in the free class, and exact finite-window detection in a structurally
controlled stress family.

\begin{thebibliography}{99}

\bibitem{DDMS1999}
J.~Dixon, M.~du~Sautoy, A.~Mann, and D.~Segal,
\emph{Analytic Pro-$p$ Groups}, 2nd ed.,
Cambridge Studies in Advanced Mathematics 61,
Cambridge University Press, 1999. doi:\href{https://doi.org/10.1017/CBO9780511470882}{10.1017/CBO9780511470882}.

\bibitem{Quadrelli2021}
C.~Quadrelli,
One-relator maximal pro-$p$ Galois groups and the Koszulity conjectures,
\emph{The Quarterly Journal of Mathematics} 72 (2021), 835--854. doi:\href{https://doi.org/10.1093/qmath/haaa049}{10.1093/qmath/haaa049}.

\bibitem{Quadrelli2022}
C.~Quadrelli,
Two families of pro-$p$ groups that are not absolute Galois groups,
\emph{Journal of Group Theory} 25 (2022), 25--62. doi:\href{https://doi.org/10.1515/jgth-2020-0186}{10.1515/jgth-2020-0186}.

\bibitem{Hamza2023}
O.~Hamza,
Zassenhaus and lower central filtrations of pro-$p$ groups considered as modules,
\emph{Journal of Algebra} 633 (2023), 172--204. doi:\href{https://doi.org/10.1016/j.jalgebra.2023.06.016}{10.1016/j.jalgebra.2023.06.016}.

\end{thebibliography}

\section{Reproducibility and verification record}

The mathematical status of the final Paper~4 core is **PASS/CLOSED**.
The repository research records retain the superseded augmentation-ideal
equality route as FAIL/CLOSED and the old order-jump route as
FAIL/CLOSED/SUPERSEDED.  The published argument uses neither.

The universal subgroup-depth comparison is proved by the weighted
normal-form filtration of the completed group algebra.  The
every-$n$ sharpness argument uses explicit Schreier-generator
commutators and a free-associative leading-word check; no abstract
graded-lift injectivity is assumed.  The intrinsic critical-window
separation uses the audited transfer predicate and its truncation
bound.

A final publication artifact should still undergo the ordinary source,
citation, LaTeX build, PDF text, page-count, reference, and checksum
audit.  These are artifact checks, not reopened mathematical gates.

\end{document}
